%% General definitions
%\documentclass{article} %% Determines the general format.
\documentclass[10pt,a4paper]{article}
\usepackage{a4wide} %% paper size: A4.
\usepackage[utf8]{inputenc} %% This file is written in UTF-8.
%% Some editors on Windows cannot save files in UTF-8.
%% If there is a problem with special characters not showing up
%% correctly, try switching "utf8" to "latin1" (ISO 8859-1).
\usepackage[T1]{fontenc} %% Format of hte resulting PDF file.
\usepackage{fancyhdr} %% Package to create a header on each page.
\usepackage{lastpage} %% Used for "Page X of Y" in the header.
%% For this to work, you have to call pdflatex twice.
\usepackage{enumerate} %% Used to change the style of enumerations (see below).

\usepackage{amssymb} %% Definitions for math symbols.
\usepackage{amsmath} %% Definitions for math symbols.
\usepackage{amsthm}
\usepackage{amsfonts}
\usepackage{mathtools}

\usepackage[english]{babel}

\usepackage{graphicx}
\usepackage{hyperref}
\setlength{\parindent}{0cm}

%\usepackage[left=2.00cm, right=2.00cm, top=2.00cm, bottom=2.00cm ]{geometry}

\newtheorem*{zz*}{z.Z}
\newtheorem*{theorem*}{Theorem}
\newtheorem*{claim*}{claim}
\newtheorem{claim}{claim}
\newtheorem*{prop*}{Propsition}
\newtheorem*{lemma*}{Lemma}
\newtheorem{lemma}{Lemma}

%\title{Lineare Algebra Serie 7}
\date{\today}


\usepackage{tikz}  %% Pagacke to create graphics (graphs, automata, etc.)
\usepackage{tikz-cd} %%commutativ system
\usetikzlibrary{automata} %% Tikz library to draw automata
\usetikzlibrary{arrows}   %% Tikz library for nicer arrow heads


%% Left side of header
\lhead{\course\\\semester\\Serie \homeworkNumber}
%% Right side of header
\rhead{\authorname\\Page \thepage\ of \pageref{LastPage}}
%% Height of header
\usepackage[headheight=36pt]{geometry}
%% Page style that uses the header
\pagestyle{fancy}

\newcommand{\authorname}{Ken Walzer}
\newcommand{\semester}{HS 2023}
\newcommand{\course}{Algebraic number Theory}
\newcommand{\homeworkNumber}{4}

%%%%%%%%%%%%%%%%%%%%%%%%%%%%%%%%%%%%%%%%5
\newcommand{\Vn}{V_{\mathbb{P}^n}}
\newcommand{\RR}{\mathbb{R}}
\newcommand{\real}{\mathrm{Re}}
\newcommand{\imag}{\mathrm{Im}}
\newcommand{\CC}{\mathbb{C}}
\newcommand{\CCC}{\mathcal{C}}
\newcommand{\QQ}{\mathbb{Q}}
\newcommand{\ZZ}{\mathbb{Z}}
\newcommand{\NN}{\mathbb{N}}
\newcommand{\FF}{\mathbb{F}}
\newcommand{\Mat}{\mathrm{Mat}}
\newcommand{\Aut}{\mathrm{Aut}}
\newcommand{\PP}{\mathbb{P}}
\newcommand{\A}{\mathbb{A}}
\newcommand{\OO}{\mathcal{O}}
\newcommand{\kk}{\mathbf{k}}
\newcommand{\aId}{\mathfrak{a}}
\newcommand{\ii}{\sqrt{-1}}
\newcommand{\GL}{\mathrm{GL}}
\newcommand{\SL}{\mathrm{SL}}
\newcommand{\orb}{\mathrm{orb}}
\newcommand{\stab}{\mathrm{stab}}
\newcommand{\ord}{\mathrm{ord}}
\newcommand{\id}{\mathrm{id}}
\linespread{1.2}

\begin{document}
	
\section*{Aufgabe 1}
	Let $F = \FF_2(T)$ and $K = F(T^{\frac{1}{2}}).$
	\begin{zz*}
		The following holds
		\begin{itemize}
			\item [(i)] $K/F$ is a field extension of degree $2.$
			\item [(ii)] $Tr_{K/F}(x) = 0$ for any $x\in K.$
			\item [(iii)] $\Delta_{K/F}(x_1, x_2) = 0$ for all $x_1, x_2\in K.$
		\end{itemize}
	\end{zz*}
	\begin{proof}
		Inhalt...
	\end{proof}
\pagebreak
%%%%%%%%%%%%%%%%%%%%%%%%%%%%%%%%%%%%%%%%%%%%%%%%%%%%%%%%%%%%
\section*{Aufgabe 2}
	Let $K$ be a number field and let $x\in \ZZ_K.$
	\begin{zz*}
		\[x\in \ZZ_K^*\Leftrightarrow N_{K/\QQ}(x) = \pm1.\]
	\end{zz*}
\pagebreak

\section*{Aufgabe 3}
	Let $m\geq 2$ be a square-free number with $m\equiv 2,3 \pmod{4}$ and $K = \QQ(\sqrt{m}).$	
	
	Let $x\in \ZZ_K^*\setminus{-1, +1}.$
	\begin{zz*}
		There exist $a,b\in \ZZ$ such that $b\neq0$ and
		\[a^2 - mb^2 = 1.\]
	\end{zz*}
\pagebreak

\section*{Aufgabe 4}
	Let $K, L$ be number fields and let $\varphi: K\to L$ a ringisomorphism.
	
	\begin{zz*}
		It holds that
		\begin{itemize}
			\item [(i)] $Tr_{K/\QQ}(x) = Tr_{L/\QQ}(\varphi(x))$ for each $x\in K.$
			\item [(ii)] $\Delta_K = \Delta_L$
		\end{itemize}
	\end{zz*}
	\begin{proof}
		Inhalt...
	\end{proof}
	
\section*{Aufgabe 5}
	Let $m,n\in \ZZ\setminus\{0,1\}$ be squarefree integers.
	\begin{zz*}
		If $\QQ(\sqrt{m}), \QQ(\sqrt{n})$ are isomorphic as rings, then $m = n.$
	\end{zz*}
	\begin{proof}
		Inhalt...
	\end{proof}
\end{document}
